\chapter{Streamlit mechanics}

Streamlit re-runs the whole script from top to bottom every time the user touches a widget. Almost every surprising
behaviour in this app comes from that fact. This chapter lists the four mechanisms the code uses to cope.

\section{Cached reads and \texttt{data\_stamp}}

Querying the database and recomputing analytics on every click would be slow with a long history. The module
\code{src/data/cached.py} wraps the read functions of \code{db/} with \code{st.cache_data}. The cache key
includes \code{db.data_stamp()}, which changes whenever the database file is written, so anything saved shows up on the
next run without manual invalidation.

\begin{lstlisting}
# src/data/cached.py
get_user_rides = _cached('get_user_rides')       # wraps db.get_user_rides
get_power_records = _cached('get_power_records')
\end{lstlisting}

\textbf{Rule:} pages \emph{read} through \code{cached.py}; code that \emph{writes} calls \code{db} directly. When you add
a new read function to \code{db/} that a page will call, add one line to \code{cached.py}.

Analytics that depend only on their inputs (the performance management chart, critical-power history) are cached on
those inputs instead, as small functions at the bottom of the same file.

\section{Sections instead of tabs}

\texttt{st.tabs()} runs the code of \emph{every} tab on every interaction and sends every tab's charts to the
browser. Pages here use a horizontal selector and an \texttt{if}, so only the visible section is built:

\begin{lstlisting}
from src.viz.sections import section_picker, remembered

sections = ["Load", "Power profile", "My new section"]
section = section_picker(sections, key="training_section")
if section == sections[0]:
    _render_load(...)
elif section == sections[2]:
    _render_my_section(...)
\end{lstlisting}

Streamlit forgets the value of a widget that is not drawn in a run. A widget inside a section must therefore be
created with \code{remembered()} so its value survives while another section is showing:

\begin{lstlisting}
period = remembered(st.selectbox, "my_period", "90 days", "Period", ["30 days", "90 days", "1 year"])
\end{lstlisting}

Do not pass \texttt{index=} or \texttt{value=} to the widget; the starting value is the third argument of
\texttt{remembered}.

The ride page's sections are drawn inside an \code{st.fragment}, so switching section does not repeat the expensive
analysis above it.

\section{Session state}

\code{st.session_state} holds things that must survive a re-run: the chosen user (\code{current_user}),
the file being analysed, a downloaded PDF waiting for its button. It does \emph{not} survive an OAuth redirect (the
browser comes back in a new session), which is why sign-in tokens are stored in the database
(\code{db.create_oauth_state}) and not in session state.

\section{Charts: automatic title and legend placement}\label{sec:figstyle}

Every function named \code{plot_*} in \code{plots.py} and \code{training_plots.py} is wrapped, at the end of the
module, by \code{figstyle.tidy_builders(globals())}. The wrapper pins the title to the top-left, puts a horizontal legend
under it and grows the top margin to fit both, so titles never overlap the plot or the legend.

\begin{itemize}
\item In those two modules, simply name your function \code{plot_something}; the tidy-up is automatic.
\item A chart built anywhere else (for example inside a page) must be wrapped by hand:
  \code{st.plotly_chart(tidy(fig))}, importing \code{tidy} from \code{src.viz.figstyle}.
\item Do not set a manual top margin or title position; they are overwritten.
\end{itemize}

\section{The look}\label{sec:theme}

\code{theme.py} holds three looks (Compact, Dashboard, Classic) as small CSS packs layered on Streamlit's theme, with
colours derived from the text colour so they work in light and dark mode. The accent colour \texttt{ACCENT} matches
\texttt{primaryColor} in \texttt{.streamlit/config.toml}; change both together. The welcome page, the FTP/LTHR help
texts and the default thresholds are in the same file.

Chart colours are a separate palette in \code{training_plots.py} (\texttt{PALETTES}), with its own light and dark
steps; plots read the current mode and call \texttt{palette(mode)}.
